% Results after the first two sections. Numbers come from numbers.tex (generated).

\subsection{Agreement with human ratings}

I compared scores with human ratings on six corpora: sentences from six genres (EmoBank\cite{buechel2017}), tweets (SemEval-2018\cite{mohammad2018semeval}), Facebook posts\cite{preotiuc2016}, film-review sentences (the Stanford Sentiment Treebank, SST\cite{socher2013}), whole film reviews with the critic's rating\cite{pang2005}, and Yelp reviews with star ratings\cite{zhang2015}. The first five carry continuous ratings and form the primary set. Every constant was fitted leave-one-corpus-out: for each target corpus, $a$, $\kappa$ and the ten coefficients came from the other five, and scores were then computed on the target's held-out test portion (Methods). The six fits chose $a$ of \locoA{} and $\kappa$ of \locoKappa.

Averaged over the five primary corpora, the calibrated score correlates with human ratings at $r=\meanSgw$, against \meanVader{} for VADER\cite{hutto2014}, \meanBlob{} for TextBlob and \meanHedo{} for the hedonometer\cite{dodds2011} (\figref{fig:accuracy}a; all values with intervals in Extended Data Table~\ref{tab:accuracy}). The plan set one criterion, which an earlier version of the method had failed: a mean $r$ at least that of VADER (Methods). It is met, \meanSgwThree{} against \meanVaderThree, by a margin that is within noise. The mean difference is \dVaderMean{} with a $t$ interval over the five corpora of \dVaderMeanCI, and on the plan's secondary outcome, Spearman's $\rho$, the two are level (\rhoSgw{} and \rhoVader; Extended Data Table~\ref{tab:spearman}). Corpus by corpus, the paired difference from VADER excludes zero in favour of the signature in \winsVader{} of six corpora and in favour of VADER in \lossesVader. On whole film reviews, the longest rated texts, TextBlob does better than both ($r=\rBlobFilm$ against \rSgwFilm{} for the signature; difference \dBlobFilm). The level of equation~(\ref{eq:level}), whose only fitted quantities are its own $a$ and $\kappa$ (\locoLevelA{} and \locoLevelKappa), reaches \meanLevel, below VADER.

Methods given more labelled text do better. A bag-of-words ridge regression trained on the same five-corpus pool as the signature's constants averages \meanBowPool, and trained on each corpus's own training portion \meanTfidf. DistilBERT fine-tuned on SST-2\cite{sanh2019} averages \meanDistil{} and a RoBERTa model fine-tuned on tweets\cite{loureiro2022} \meanRoberta. The lead of the transformers is uneven. The signature is ahead of DistilBERT on tweets and on Facebook posts (differences \dDistilTw{} and \dDistilFb), level with it on EmoBank, and far behind on SST, the domain DistilBERT was tuned on, and on reviews. The signature is not a replacement for these models where accuracy on one text matters most.

Removing parts of the signature shows what earns its place (\figref{fig:accuracy}c; means over all six corpora). Without the operator channels the mean correlation falls by \lossNoShift, and without dilution ($\kappa=0$) by \lossNoDil. Without the neutral prior ($a=0$) it falls by \lossNoShrink, a mean over a gain from the prior in \winsNoShrink{} corpora, a small loss in film reviews and no clear change in the other three. Merging hedons and dolors into one channel costs \lossOneChan{} and dropping the position moments \lossNoPos: neither matters for agreement with raters, though both are needed for the readouts below.

The fitted coefficients can be read as rules, with a caution: the readings belong to one pooled fit, the shipped model, and the bootstrap intervals describe sampling error within that fit and say nothing about other corpora. In it a dolor moves the predicted rating \constDolor{} \constDolorCI{} times as far as a hedon. A negated positive word counts as \constNegH{} \constNegHCI{} times its charge and a negated negative word as \constNegD{} \constNegDCI: \emph{not good} is bad, and \emph{not bad} is neutral. An intensifier adds \constBoost{} \constBoostCI{} times the charge of the word it modifies. Across the six leave-one-corpus-out fits the first three readings run from \locoDolorLo{} to \locoDolorHi, from \locoNegHLo{} to \locoNegHHi{} and from \locoNegDLo{} to \locoNegDHi, and fits to single corpora differ far more (Extended Data Table~\ref{tab:readings}): the ratio of dolor to hedon is below 1 in SST and above 4 in Yelp. One reading has the same sign in every fit, pooled or single: a negated positive word turns negative. The sizes are properties of these corpora. The data do not settle the effect of diminishers (\constDamp{} \constDampCI) or of contrast words (\constPost{} \constPostCI); in particular they give no support to the rule that the clause after \emph{but} counts for more\cite{hutto2014}.

Word charges can also be fitted. With the exposure and the operator constants held fixed, the predicted rating is linear in the vector of charges, so the charges that best predict ratings, pulled toward the human word norms by a ridge penalty, solve one convex problem (Methods). Fitted leave-one-corpus-out, these refined charges raise the mean correlation to \meanRefined, the same as the bag-of-words ridge trained on the same five corpora (\meanBowPool). The gain is uneven. It is large on review text (SST \rSgwSst{} to \rRefinedSst, film reviews \rSgwFilm{} to \rRefinedFilm) and negative on the mixed genres of EmoBank (\rSgwEmo{} to \rRefinedEmo) and on Facebook posts (\rSgwFb{} to \rRefinedFb). Charges learned from ratings carry the genre of the texts they were learned from, and their attitudes. Of \grpTerms{} words for groups of people and two public figures, checked after a referee agent raised the question, \grpDown{} moved down by more than 0.05 on the charge scale (\emph{women} from \grpWomenNorm{} to \grpWomenFit, \emph{muslim} from \grpMuslimNorm{} to \grpMuslimFit; Extended Data Table~\ref{tab:groups}). Human word norms travel better, and the refined charges are not part of the released package.

One hoped-for property did not hold. The standard error of a level shrinks with exposure, and the error of the calibrated score against human raters does not follow it. From the lowest to the highest quartile of exposure the error falls in \errFall{} corpora (\errFallNames) and rises in the other \errRise{} (Extended Data Fig.~\ref{fig:extended}a). The standard error says how well a text's own words pin down its level. It is not a prediction interval for a human rating.

\figAccuracy

\subsection{Cost}

Scoring is one table lookup per word. On one core of a rented six-core desktop processor the implementation, written in Python with NumPy, turns \sgwRate{} sentences per second, or \sgwWords{} million words, into signatures and levels (\figref{fig:accuracy}b). On the same core VADER scores \vaderRate{} sentences per second, TextBlob \blobRate, DistilBERT \distilRate{} and the RoBERTa model \robertaRate. The signature is therefore about \speedDistil{} times faster than DistilBERT and \speedRoberta{} times faster than the RoBERTa model, well past the factor of 100 that the plan set as its criterion. On a graphics card (\gpuName, half precision) DistilBERT reaches \distilGpuHalf{} sentences per second, and one processor core running the signature is still \speedDistilGpu{} times faster. Time grows in proportion to the amount of text (log-log slope \scalingSlope{} from $10^4$ to $10^6$ sentences); ten million sentences took \streamSeconds{} seconds in fixed-size batches. Because signatures merge, the work divides among processes without changing any result, and twelve processes on the six cores scored \sgwAllCores{} million sentences per second, \parSpeedup{} times the rate of one. At these rates a billion sentences take \billionHours{} hours on one core or \billionMinutesAll{} minutes on the six, against \billionHoursDistilGpu{} hours for DistilBERT on the graphics card and \billionDaysDistil{} days on one core.

\subsection{Where the feeling sits}

Front-loading makes three predictions, and I tested each (\figref{fig:position}).

The first is that charged words come early. They do not. The charge-weighted mean position, on a scale from $-1$ at the first word to $+1$ at the last, is above zero in all \whereUnits{} kinds of unit examined (the six rated corpora, paragraphs of the film reviews and chapters of 150 of the novels), from \whereMin{} in Yelp reviews to \whereMax{} in tweets, with every interval excluding zero (Extended Data Table~\ref{tab:where}). In EmoBank sentences, SST sentences and tweets the first tenth of the words carries \firstDecileLo{} to \firstDecileHi{} of the text's mean charge per word, and in those three and in Facebook posts the last tenth carries \lastDecileLo{} to \lastDecileHi{} (\figref{fig:position}a). English sentences tend to open with determiners and pronouns and close on content. Facebook posts open at the mean (\fbFirst). Yelp reviews are the one place with a charged opening, at \yelpFirst{} times the mean, matched by the close at \yelpLast. Negative charge sits later than positive charge in \whereDolorLater{} of the \whereUnits.

The second prediction is that readers weigh early words more. Here the answer depends on the kind of text. The position kernel fitted in each corpus puts more weight on the first third of a text than on the last third in SST sentences (difference \kernSst{} \kernSstCI), as front-loading predicts. It puts less in tweets (\kernTw{} \kernTwCI) and in Yelp reviews (\kernYelp{} \kernYelpCI), and is undecided in the other three (\figref{fig:position}c). In the shipped model, fitted on all six corpora, the last third outweighs the first (\constPos{} \constPosCI).

A third analysis, added after the first results, uses human ratings alone. The SST gives a human rating to every phrase of its parse trees\cite{socher2013}, so the rating of a phrase can be regressed on the ratings of its two parts. In the training trees the later part gets the larger weight, \wRightTwo{} against \wLeftTwo{} where both parts have at least two words (\ptNTrain{} phrase pairs). The comparison is confounded by length. The later part is the longer one in \ptLaterLonger\% of pairs (\ptWordsLater{} words on average against \ptWordsEarlier), and the weight of a part rises with its share of the phrase's words, from \ptShareLo{} for a part with under a tenth of them to \ptShareHi{} for a part with over nine tenths (\figref{fig:position}b; all trees, \ptN{} pairs). Where the two parts are equally long, the earlier one has the larger weight: \ptEqEarlier{} against \ptEqLater{} (later minus earlier, \ptEqDiff{} \ptEqDiffCI; \ptEqN{} pairs). With one curve of weight against word share for both parts, the extra weight of the later part is \ptOrder{} \ptOrderCI. So in film-review phrases the longer part counts for more, and length for length the earlier part counts for slightly more. English syntax may be the cause as much as order, since the two parts of a phrase differ in grammatical role. The dependence on length is modest: a plain mean of the two part ratings accounts for the phrase rating as well as a mean weighted by length ($R^2=\ptRsqPlain$ and \ptRsqLength; \ptRsqFree{} with free weights).

The plan's third test, within EmoBank documents, is undecided: rated intensity neither falls nor rises with sentence position (slope \emoSlope{} \emoSlopeCI{} per document length; \emoDocs{} documents).

Front-loading is therefore refuted as a claim about where writers put feeling words. As a claim about what readers weigh it holds in film-review sentences, by the kernel and by the phrase ratings at equal length, fails in tweets and in Yelp reviews, and is undecided elsewhere. For whole reviews, earlier work found that the closing sentences carry the verdict\cite{pang2004,taboada2004,raychev2009,becker2010}.

\figPosition

\subsection{Neutral words dilute, from either side}

Dilution fares better (\figref{fig:dilution}). The planned test asks whether a readout with a buffer weight $\kappa>0$ agrees with held-out ratings better than the same readout with $\kappa=0$. It does in \winsNoDil{} of six corpora (tweets, Facebook posts, SST and Yelp) and is undecided in EmoBank (\dNoDilEmo) and in film reviews (\dNoDilFilm). Two statistics that the plan did not list agree in direction. Holding a text's total charge fixed, rated intensity falls as the number of buffer words rises in \partialNeg{} of six corpora, EmoBank among them, with partial rank correlations from \partialEmo{} in EmoBank to \partialSst{} in SST sentences (\figref{fig:dilution}b); whole film reviews show no effect (\partialFilm{} \partialFilmCI). In fits to single corpora, agreement is highest at a buffer weight between \kappaLo{} and \kappaHi{} in the same five and at zero for film reviews, with gains in $r$ of \gainKappaLo{} to \gainKappaHi{} that carry no intervals (\figref{fig:dilution}c).

Human phrase ratings show the effect without any scorer. Take one rated word of the treebank joined to a neutral phrase of $m$ words. With the neutral words after it, the joined phrase keeps \ptKeptAfterFirst{} \ptKeptAfterFirstCI{} of the word's rating for $m=1$ and \ptKeptAfterLast{} \ptKeptAfterLastCI{} for $m$ of \ptKeptAfterLastM. With the neutral words before it, the share is \ptKeptBeforeFirst{} \ptKeptBeforeFirstCI{} for $m=1$ and \ptKeptBeforeLast{} \ptKeptBeforeLastCI{} for $m$ of \ptKeptBeforeLastM{} (\figref{fig:dilution}a). The share falls with $m$ on both sides. Over rated parts of every length, matched on the lengths of both parts, the share kept is \ptDil{} \ptDilCI{} lower when the neutral words come first than when they follow. The level of these shares is pulled down by regression to the mean, because the rated part was selected for a strong rating; the trend with $m$ and the comparison between sides are not.

The 2024 manuscript had dilution right. Words that carry no feeling weaken the judged feeling of the whole, more of them weaken it more, and it does not take a particular order, as the averaging account of impression formation predicts\cite{anderson1965,nisbett1981}.

\subsection{Recency, proposed and withdrawn}

Two of the first results, charged words sitting late and the later part of a treebank phrase carrying the larger weight, suggested a hypothesis opposite to front-loading: recency, that what comes last in a text counts most toward a reader's judgement of it. I wrote it down with three predictions and committed them before examining the data they concern, the development and test trees of the treebank and \confReviews{} IMDB film reviews (Methods). By their stated rules all three held. In \confNodes{} held-out phrase pairs the later part has weight \confLater{} and the earlier part \confEarlier{} (R1; difference \confWeight{} \confWeightCI). A rated phrase keeps \confKeptAfter{} of its rating when two or more neutral words follow it and \confKeptBefore{} when they come before it (R2; difference \confKept{} \confKeptCI). And the last fifth of an IMDB review separates favourable from unfavourable reviews better than the first fifth does (R3; area under the curve \confAucLast{} against \confAucFirst; difference \confAuc{} \confAucCI).

A referee agent then pointed out that R1 and R2 do not hold length fixed. With length fixed, both fail. In the held-out trees the later part is the longer one in \ptLaterLongerHeld\% of pairs. At equal length the difference in R1 reverses (\ptEqDiffHeld{} \ptEqDiffHeldCI), and with a common curve of weight against word share the later part's extra weight is \ptOrderHeld{} \ptOrderHeldCI. In R2 the rated parts that were followed by neutral words averaged \ptRatedAfter{} words and those preceded by them \ptRatedBefore, and a short rated part keeps less of its rating whichever side the neutral words are on. Matched on both lengths the difference reverses (before minus after, \ptDilHeld{} \ptDilHeldCI).

R3 compares fifths of equal length and stands. It measures where reviewers place their verdict, which is a fact about writers. The scorer's density in the first, the three middle and the last fifth separates the two classes with areas of \imdbFifthOne, \imdbMiddleLo{} to \imdbMiddleHi{} and \imdbFifthFive: the verdict shows at both ends of a review and most at the end, as the earlier work on reviews found\cite{pang2004,becker2010}.

Recency as a claim about readers is withdrawn. Its tests were written down in advance and passed, and they were the wrong tests, because the plan left out the control they needed. Both steps are on record in the analysis plan.

\figDilution

\subsection{A book at every scale}

I scored \nNovels{} novels from Project Gutenberg (\novelWords{} million words, median \novelChapters{} chapters) sentence by sentence and merged the signatures up to chapters and books. \figref{fig:book}a,b shows \emph{Moby-Dick}: \mobyWords{} words carrying \mobyHedons\hd{} and \mobyDolors\dl, a density of \mobyDensity{} per thousand words, with each of the \mobyChapters{} sections that the chapter rule finds in it (the book has 135 chapters and an epilogue; Methods) and twelve spans of the last one read from the same pass.

Chapters are not random samples of a book's sentences. The variance of chapter densities is \heteroMedian{} times (median; interquartile range \heteroLo{} to \heteroHi) what results when a novel's sentences are shuffled across its chapters, and exceeds the shuffled value at $p\leq0.01$ in \heteroSig\% of novels (\figref{fig:book}c). The shuffle destroys all continuity between neighbouring sentences, so the ratio shows that tone is locally coherent in novels. It does not show that chapter boundaries are where tone changes, which would need a null that moves the boundaries and keeps the order.

The cubic arc read from a book's signature captures \arcShareThree\% (median) of the variance of an arc drawn with a moving window one tenth of the book wide; keeping moments up to degree nine would capture \arcShareNine\% (\figref{fig:book}d). The two leading modes of variation among the window arcs of the \nNovels{} novels carry \svdVarOne\% and \svdVarSecond\% of the variance and are close to the first two Legendre polynomials (cosine similarity \svdCosOne{} and \svdCosTwo), the rise or fall and the single turn that earlier studies of story shapes found\cite{reagan2016}; the third mode is not close to the cubic (\svdCosThree). A wide moving window yields smooth leading modes from any series, so the resemblance is weak evidence about stories.

Human labels exist for one body of narrative. Alm and colleagues labelled every sentence of \nTales{} fairy tales with an emotion\cite{alm2005}. The cubic arc of each tale's merged signature agrees in sign with the cubic arc of its human labels in \taleArcPos\% of tales (median correlation \taleArc; sign test $p=$~\taleArcP; Extended Data Fig.~\ref{fig:extended}b). For a tale's overall valence the picture is less favourable. The calibrated score of the merged signature correlates with the mean human label at \taleStoryCal{} \taleStoryCalCI, and VADER averaged over sentences does better (\taleStoryVader{} \taleStoryVaderCI). The human target here is itself a mean of sentence labels, which an average of sentence scores mirrors. VADER applied to the whole tale saturates beyond $\pm0.99$ for \taleSaturated\% of tales and falls to \taleWholeVader.

A caution applies to every number in this section. Word norms rate words, not events. The final chapter of \emph{Moby-Dick}, in which the ship is destroyed, scores near the book's own level, because \emph{whale}, \emph{sea} and \emph{boat} are pleasant words to raters. Densities of narrative text are best read against the text's own baseline.

\figBook

\subsection{Pooling, averaging and chunking}

Bounded scores depend on how a long text is cut (\figref{fig:aggregation}a to c). For the first \vaderBooks{} novels in file order I computed a book score as the mean VADER compound score over chunks. The mean across books rises from \vaderOne{} with one-sentence chunks to \vaderChapter{} with whole chapters, and \vaderFlipChapter\% of pairs of books swap order between the two. The same happens on \partsReviews{} whole film reviews. The score of a review and the mean score of its sentences differ in sign for \partsDistilSign\% of reviews with DistilBERT, \partsRobertaSign\% with the RoBERTa model and \partsVaderSign\% with VADER, whose whole-review score is beyond $\pm0.99$ for \partsVaderSaturated\% of them. For the transformers the comparison is confounded by truncation: \truncShare\% of the reviews exceed the models' 512-token window, so the whole-review score reads the opening of a review and the sentence mean reads all of it. In all three cases the mean of sentence scores predicts the critic's rating better than the whole-review score does ($r=\partsDistilSent$ against \partsDistilWhole, \partsRobertaSent{} against \partsRobertaWhole, and \partsVaderSent{} against \partsVaderWhole).

What depends on the cut is the average. If the densities of chunks are averaged, \densFlipLast\% of pairs of books swap order between one-sentence chunks and chunks of 125 sentences, against \vaderFlipLast\% for VADER (\densFlipChapter\% and \vaderFlipChapter\% with whole chapters). If the chunks' signatures are merged, the book's density is the same for every chunking, to a relative difference below \sgwSpread.

Invariance is not validity. On the same \partsReviews{} reviews the merged density correlates with the critic's rating at \partsSgwDensity, far below the mean of DistilBERT's sentence scores (\partsDistilSent), and among the fairy tales an average of VADER's sentence scores tracked overall valence better than the merged signature did. Merging removes the dependence on the cut. It does not make a score built from word ratings a better reader.

How far apart are a pooled value and a mean of parts? For parts with exposures $D_p$ and densities $\lambda_p$,
\begin{equation}
\frac{\sum_p D_p\lambda_p}{\sum_p D_p}=\overline{\lambda}+\frac{\Cov(D,\lambda)}{\overline{D}},
\label{eq:pooled}
\end{equation}
an identity familiar from the decomposition of aggregate productivity\cite{olley1996}. The two differ when long parts differ in density from short ones. In novels they rarely do: the median gap between a novel's pooled density and the mean of its chapter densities is \gapOverSd\% of the standard deviation of density between novels, and \pairFlip\% of the \nPairs{} pairs of novels are ordered differently by the two (\figref{fig:aggregation}d).

A difference between two bodies of text splits in the same spirit. Kitagawa's decomposition\cite{kitagawa1955} writes a difference in pooled density as a part due to differences within strata plus a part due to the mix of strata. Four film critics in the review corpus\cite{pang2005} illustrate it, with the critic's own rating as the stratum (\figref{fig:aggregation}e). Pooled, Berardinelli's and Renshaw's reviews have the same density (difference \critBerRenTotal{} per thousand words). Within the same rating Berardinelli's prose is the more negative by \critBerRenWithin, an interval that reaches zero, and his mix of ratings adds back \critBerRenMix. None of the \critPairs{} pairs of critics reverses between the pooled and the within-rating comparison. The decomposition is bookkeeping, and which part answers a question depends on the question\cite{pearl2014}; a signature kept per stratum makes both parts available, and a pooled bounded score does not.


\figAggregation
